\inicializarSubseccion{Conclusión de Alexa}

\tab En este reto aprendí el proceso de obtener el volumen de un sólido de revolución a través de código, que para llegar a eso pase por conceptos que tampoco conocía como la interpolación y funciones que no conocía, tanto en matlab como en geogebra, por ejemplo, en matlab no conocía la función vander (para sacar los coeficientes) y cylinder (para graficar) pero sobre todo no sabía sobre el proceso que hicimos en geogebra, ya que había utilizado este software antes pero solo para graficar funciones por separado. No sabía que se podían crear puntos para después unirlos con la función polynomial. Entonces más que nada aprendí nuevas herramientas de estos dos softwares. Ahora, sobre el proceso para encontrar el volumen (el área bajo la curva) es algo que ya conocía pero nunca lo había aplicado a código.

\tab Por esto mismo de que no había aplicado este conocimiento en código y que no conocía bien las herramientas que usamos, mis principales dificultades fueron errores en el código pero errores que eran más por accidentes que por lógica, por ejemplo escribir x1 en vez de y1, esto hacía que mis resultados se modificarán, y como eran errores tan pequeños, revisaba el código y no los veía, entonces pensaba que eran errores de lógica y cambiaba cosas que no tenía que modificar y se hacía un caos, lo que me llevaba mucho tiempo.

\tab En conclusión mis principales dificultades fueron accidentes en el código que me hacían pensar que tenía errores de lógica, pero fuera de eso los talleres y el equipo lo hicieron mucho más comprensible y fácil.

\tab Por último, los cálculos al final no nos dieron mucho problema, a excepción del costo del material, el cual hizo que el precio final fuera bastante elevado, por lo que influyó en la decisión final del material. Fuera de esto, todos los demás cálculos fueron bastante razonables y reales. 
